% Options for packages loaded elsewhere
\PassOptionsToPackage{unicode}{hyperref}
\PassOptionsToPackage{hyphens}{url}
\documentclass[
  ignorenonframetext,
]{beamer}
\newif\ifbibliography
\usepackage{pgfpages}
\setbeamertemplate{caption}[numbered]
\setbeamertemplate{caption label separator}{: }
\setbeamercolor{caption name}{fg=normal text.fg}
\beamertemplatenavigationsymbolsempty
% remove section numbering
\setbeamertemplate{part page}{
  \centering
  \begin{beamercolorbox}[sep=16pt,center]{part title}
    \usebeamerfont{part title}\insertpart\par
  \end{beamercolorbox}
}
\setbeamertemplate{section page}{
  \centering
  \begin{beamercolorbox}[sep=12pt,center]{section title}
    \usebeamerfont{section title}\insertsection\par
  \end{beamercolorbox}
}
\setbeamertemplate{subsection page}{
  \centering
  \begin{beamercolorbox}[sep=8pt,center]{subsection title}
    \usebeamerfont{subsection title}\insertsubsection\par
  \end{beamercolorbox}
}
% Prevent slide breaks in the middle of a paragraph
\widowpenalties 1 10000
\raggedbottom
\AtBeginPart{
  \frame{\partpage}
}
\AtBeginSection{
  \ifbibliography
  \else
    \frame{\sectionpage}
  \fi
}
\AtBeginSubsection{
  \frame{\subsectionpage}
}
\usepackage{iftex}
\ifPDFTeX
  \usepackage[T1]{fontenc}
  \usepackage[utf8]{inputenc}
  \usepackage{textcomp} % provide euro and other symbols
\else % if luatex or xetex
  \usepackage{unicode-math} % this also loads fontspec
  \defaultfontfeatures{Scale=MatchLowercase}
  \defaultfontfeatures[\rmfamily]{Ligatures=TeX,Scale=1}
\fi
\usepackage{lmodern}
\ifPDFTeX\else
  % xetex/luatex font selection
\fi
% Use upquote if available, for straight quotes in verbatim environments
\IfFileExists{upquote.sty}{\usepackage{upquote}}{}
\IfFileExists{microtype.sty}{% use microtype if available
  \usepackage[]{microtype}
  \UseMicrotypeSet[protrusion]{basicmath} % disable protrusion for tt fonts
}{}
\makeatletter
\@ifundefined{KOMAClassName}{% if non-KOMA class
  \IfFileExists{parskip.sty}{%
    \usepackage{parskip}
  }{% else
    \setlength{\parindent}{0pt}
    \setlength{\parskip}{6pt plus 2pt minus 1pt}}
}{% if KOMA class
  \KOMAoptions{parskip=half}}
\makeatother
\setlength{\emergencystretch}{3em} % prevent overfull lines
\providecommand{\tightlist}{%
  \setlength{\itemsep}{0pt}\setlength{\parskip}{0pt}}
\usepackage{bookmark}
\IfFileExists{xurl.sty}{\usepackage{xurl}}{} % add URL line breaks if available
\urlstyle{same}
\hypersetup{
  hidelinks,
  pdfcreator={LaTeX via pandoc}}

\author{\texorpdfstring{}{}}
\date{}

\begin{document}

\begin{frame}[fragile]{\(q^2\) error scale (Lagrange / Hurwitz
viewpoint)}
\protect\phantomsection\label{q2-error-scale-lagrange-hurwitz-viewpoint}
For integers \(p,q\) with \(q \geq 1\), define the \textbf{weighted}
error \begin{equation}
\label{eq:lagrange}
\mu_Q(x) = \min_{1 \leq q \leq Q}\ \min_{p \in \mathbb{Z}} q^2 \left| x - \frac{p}{q} \right|.
\end{equation}

For fixed \(q\), the inner minimum is attained at
\(p = \mathrm{round}(qx)\), so the implementation checks the same three
candidates per \(q\) as for \(\delta_Q(x)\) in
\texttt{02\_problem\_statement.md}. The factor \(q^2\) makes \(\mu_Q\)
dimensionless in the sense used in Hurwitz-type theorems: for infinitely
many convergents one has \(q^2|x - p/q| < 1\), and for every irrational
\(x\) there are infinitely many pairs with \(q^2|x-p/q| < 1/\sqrt{5}\)
{[}@hurwitz1891{]}.

\begin{block}{Interpreting \(\mu_Q\) at finite \(Q\)}
\protect\phantomsection\label{interpreting-mu_q-at-finite-q}
At \textbf{finite} \(Q\), \(\mu_Q(x)\) answers: \emph{how small is the
best weighted error achievable with denominators up to \(Q\)?} A superb
approximant with small \(q\) can yield a \textbf{moderate} \(\delta_Q\)
but a \textbf{very small} \(\mu_Q\) because \(q^2\) penalises large
denominators less aggressively than \(1/q\) rewards them in
\(\delta_Q\). Conversely, a number may have a tiny \(\delta_Q\) because
some modest \(q\) lands unusually close, while \(\mu_Q\) need not be
extreme if that hit is not outstanding on the \(q^2\) scale.

The Monte Carlo pipeline records both \(\delta_Q\) and \(\mu_Q\) for
named constants and for the pooled reference
(\texttt{min\_q\_squared\_error},
\texttt{empirical\_percentile\_rank\_q\_squared},
\texttt{reference\_q\_squared\_summary} in
\texttt{output/data/proximity\_summary.json}). Figure
\ref{fig:proximity_mu} shows the pooled reference distribution of
\(\log_{10}\mu_Q\) with the same constant markers as the \(\delta_Q\)
panels.
\end{block}

\begin{block}{Comparison to badly approximable heuristics}
\protect\phantomsection\label{comparison-to-badly-approximable-heuristics}
Badly approximable numbers satisfy \(\liminf_q q^2|x-p/q|>0\), so
\(\mu_Q(x)\) is bounded away from \(0\) as \(Q\to\infty\) along
worst-case subsequences. At fixed moderate \(Q\), empirical ranks of
\(\mu_Q\) still fluctuate with the \textbf{best} hit available below the
ceiling; reporting both statistics reduces the risk of over-interpreting
a single scale.
\end{block}
\end{frame}

\end{document}
